\section{Symmetric tensors on split elliptic ruled surfaces}
\label{sec:split-elliptic-symmetric-tensors}

Let $E=\mathbb C/\Lambda$, where $\Lambda=\mathbb Z+\mathbb Z\tau$ and $\operatorname{Im}\tau>0$. Let $A$ have positive degree $d$, and put $M=\mathbb P_{\mathrm{lines}}(\mathcal O_E\oplus A)$. Use the normalized factors of \cref{lem:elliptic-degree-automorphy-normalization}, so that
\[
 b_\lambda(z)=\chi_\lambda
       \exp\bigl(-\pi\mathrm{i}d(m^2\tau+2mz)\bigr),
 \qquad q_\lambda=-b_\lambda'/b_\lambda=2\pi\mathrm{i}dm,
 \qquad \lambda=m\tau+n.
\]
A symmetric tangent tensor is a polynomial function on the cotangent bundle, homogeneous in its cotangent fiber variables. Write $u\,\mathrm{d}z+y\,\mathrm{d}t$ for a cotangent vector in the finite ruling chart. The Euler field $t\partial/\partial t$ gives the global linear function $h=ty$. The ruling action decomposes these polynomials into weights. Their dimensions measure the abundance of symmetric tensors, while their products detect the obstruction to finite generation.

Let $U$ be the rank-two bundle on $E$ whose sections on the universal cover are linear polynomials in formal variables $h,u$ satisfying
\[
 F(z+\lambda,h,u)=F(z,h,u-q_\lambda h).
\]
Thus the invariant polynomial $h$ and the coefficient of $u$ give an exact sequence $0\to\mathcal O_Eh\to U\to\mathcal O_E\to0$.

\begin{lemma}\label{lem:split-elliptic-tensor-weights}
For every integer $r\geq0$, the space $H^0(M,\operatorname{Sym}^rT_M)$ has a finite weight decomposition with weights $-r\leq\ell\leq r$. The weight-zero space is $H^0(E,\operatorname{Sym}^rU)$. For $1\leq k\leq r$, the weight $-k$ space is
\[
 H^0(E,A^k\otimes\operatorname{Sym}^{r-k}U),
\]
and the weight $k$ space is $H^0(E,A^{-k}\otimes\operatorname{Sym}^{r-k}U)$. Multiplication adds the tensor degrees and the weights.
\end{lemma}
\begin{proof}
On the cover $\mathbb C\times\mathbb P^1$, ordinary cotangent coordinates change by
\[
 v=1/t,\qquad u_\infty=u,\qquad y_\infty=-t^2y.
\]
The lattice acts by
\[
 (z,t,u,y)\longmapsto
 (z+\lambda,b_\lambda t,u+q_\lambda ty,y/b_\lambda).
\]
These formulas follow by preserving the one-form $u\,\mathrm{d}z+y\,\mathrm{d}t$. In the other ruling chart the action is
\[
 (z,v,u_\infty,y_\infty)\longmapsto
 (z+\lambda,b_\lambda^{-1}v,
 u_\infty-q_\lambda v y_\infty,b_\lambda y_\infty).
\]
It is holomorphic at $v=0$ and agrees with the finite-chart formula. Both formulas preserve $h=ty=-v y_\infty$.

A degree-$r$ homogeneous function is a sum of terms $g_b(z,t)y^bu^{r-b}$ for $0\leq b\leq r$, with $g_b$ holomorphic on $\mathbb C^2$. Holomorphy at $v=0$ says that $v^{2b}g_b(z,1/v)$ is holomorphic there. Its Taylor or Laurent expansion consequently gives
\[
 g_b(z,t)=\sum_{a=0}^{2b}c_{a,b}(z)t^a,
\]
with entire coefficients. No cancellation between different $b$ is possible, since the monomials in $u_\infty,y_\infty$ are independent.

Give $t$ weight $1$, $y$ weight $-1$ and $u$ weight $0$. A term $t^ay^bu^{r-b}$ has weight $\ell=a-b$. Both chart changes and the lattice preserve this grading; equivalently, it is the cotangent lift of the global action $t\mapsto at$ for $a\in\mathbb C^*$. The bounds on $a$ imply $-r\leq\ell\leq r$.

For $\ell=-k<0$, a weight block is uniquely $y^kF(z,h,u)$, with $F$ homogeneous of degree $r-k$. Its invariance is equivalent to
\[
 F(z+\lambda,h,u)
       =b_\lambda(z)^kF(z,h,u-q_\lambda h).
\]
This is the factor for $A^k\otimes\operatorname{Sym}^{r-k}U$. Conversely, each such polynomial defines a function holomorphic in both ruling charts: the exponents satisfy $0\leq a\leq2b$. For $\ell=k>0$, a weight block is uniquely $t^kh^kF(z,h,u)$ of the same remaining degree. Its invariance has the factor $b_\lambda^{-k}$, giving $A^{-k}\otimes\operatorname{Sym}^{r-k}U$. Weight zero consists precisely of homogeneous polynomials $F(z,h,u)$ with the factor defining $\operatorname{Sym}^rU$. Products of polynomial functions add degrees and weights, proving the last assertion.
\end{proof}

\begin{lemma}\label{lem:elliptic-unipotent-symmetric-cohomology}
The extension $U$ is nonsplit. For every $s\geq0$, one has
\[
 H^0(E,\operatorname{Sym}^sU)=\mathbb C h^s.
\]
For $k\geq1$ and $s\geq0$, one has
\[
 \begin{aligned}
 H^0(E,A^{-k}\otimes\operatorname{Sym}^sU)&=0,\\
 H^1(E,A^k\otimes\operatorname{Sym}^sU)&=0,\\
 h^0(E,A^k\otimes\operatorname{Sym}^sU)&=kd(s+1).
 \end{aligned}
\]
For $s\geq1$, the map taking the coefficient of $u^s$ gives a surjection
\[
 H^0(E,A\otimes\operatorname{Sym}^sU)\longrightarrow H^0(E,A),
\]
whose kernel is $hH^0(E,A\otimes\operatorname{Sym}^{s-1}U)$.
\end{lemma}
\begin{proof}
A splitting of $U$ would give a section $u+a(z)h$ lifting $1$, with $a$ entire and
\[
 a(z+1)=a(z),\qquad a(z+\tau)=a(z)-2\pi\mathrm{i}d.
\]
Differentiating shows that $a'$ is lattice periodic, hence constant on the compact torus. The first identity forces that constant to vanish, contradicting the second. Thus the extension is nonsplit.

More generally, write an invariant polynomial of degree $s$ as
\[
 F(z,h,u)=\sum_{j=0}^s a_j(z)h^{s-j}u^j.
\]
If its largest nonzero $u$-degree is $j>0$, then $a_j$ is lattice periodic and therefore a nonzero constant. Comparing the next coefficient gives
\[
 a_{j-1}(z+1)=a_{j-1}(z),\qquad
 a_{j-1}(z+\tau)=a_{j-1}(z)-j(2\pi\mathrm{i}d)a_j.
\]
The same derivative argument gives a contradiction. Hence $F=a_0h^s$, with $a_0$ constant. This also treats $s=0$.

The powers of the invariant line $\mathcal O_Eh$ give a filtration of $\operatorname{Sym}^sU$ with $s+1$ trivial line quotients. A negative-degree line bundle on a compact curve has no nonzero section, since the zero divisor of a nonzero section has nonnegative degree. The filtration twisted by $A^{-k}$ thus has no global section. Serre duality and $\Omega_E^1\simeq\mathcal O_E$ give $H^1(E,A^k)=0$. The positive twist of the filtration and its cohomology exact sequences then give $H^1(E,A^k\otimes\operatorname{Sym}^sU)=0$. Riemann--Roch on the genus-one curve, with rank $s+1$ and degree $kd(s+1)$, gives the dimension formula.

Finally, the leading coefficient gives the bundle exact sequence
\[
 0\longrightarrow A\otimes\operatorname{Sym}^{s-1}U
 \xrightarrow{\ h\ }A\otimes\operatorname{Sym}^sU
 \longrightarrow A\longrightarrow0.
\]
Indeed, translation of $u$ by a multiple of $h$ leaves its leading coefficient unchanged, and a polynomial with zero leading coefficient is divisible by $h$. The vanishing just proved gives the asserted surjectivity and kernel on global sections.
\end{proof}

\begin{theorem}\label{thm:split-elliptic-symmetric-dimensions}
For every integer $r\geq0$,
\[
 h^0(M,\operatorname{Sym}^rT_M)
       =1+\frac{d\,r(r+1)(r+2)}6.
\]
In the quotient convention for $\pi:\mathbb P(T_M)\to M$, the tautological line bundle $\mathcal O_{\mathbb P(T_M)}(1)$ is big and has volume $d$.
\end{theorem}
\begin{proof}
The positive weights vanish, the zero weight contributes one dimension and the weight $-k$ contributes $kd(r-k+1)$. Thus
\[
 h^0(M,\operatorname{Sym}^rT_M)
 =1+d\sum_{k=1}^r k(r-k+1)
 =1+\frac{d\,r(r+1)(r+2)}6.
\]
The elliptic curve and its holomorphic line bundle are algebraic, so $M$ and $\mathbb P(T_M)$ are projective. The projective-bundle identity $\pi_*\mathcal O_{\mathbb P(T_M)}(r)=\operatorname{Sym}^rT_M$ identifies their section spaces. Since $\mathbb P(T_M)$ has dimension three, the definition of volume gives
\[
 \operatorname{vol}\bigl(\mathcal O_{\mathbb P(T_M)}(1)\bigr)
 =\limsup_{r\to\infty}\frac{6h^0(M,\operatorname{Sym}^rT_M)}{r^3}=d>0.
\]
Positive volume is the section-growth criterion for bigness.
\end{proof}

\begin{theorem}\label{thm:split-elliptic-tangent-ring-nonfinite}
The complex algebra
\[
 R=\bigoplus_{r\geq0}H^0(M,\operatorname{Sym}^rT_M)
\]
is not finitely generated.
\end{theorem}
\begin{proof}
Write $R_{r,k}$ for tensor degree $r$ and weight $-k$. The surviving weights have $k\geq0$, and multiplication adds these indices. The zero-weight subalgebra is $\mathbb C[h]$. Let $I$ be the ideal consisting of all components of positive $k$, and form
\[
 S=R/(hR+I^2).
\]
Its zero-weight component is $\mathbb C$. Its positive-weight ideal $J$ satisfies $J^2=0$, so $S=\mathbb C\oplus J$ as a vector space.

For every $r\geq1$, products in $I^2$ have weight index at least two. Consequently, the degree-$r$, weight-index-one component of $S$ is
\[
 R_{r,1}/hR_{r-1,1}\simeq H^0(E,A).
\]
For $r=1$, this is the identity and $R_{0,1}=0$. For $r\geq2$, it follows from the last assertion of \cref{lem:elliptic-unipotent-symmetric-cohomology}, with $s=r-1$. Each of these components has dimension $d>0$. Since their degrees are distinct, $J$ is infinite dimensional.

If $R$ were finitely generated, its quotient $S$ would be finitely generated. Write a finite set of algebra generators of $S$ as $c_i+j_i$, with $c_i\in\mathbb C$ and $j_i\in J$. Because $J^2=0$, every polynomial in these generators lies in the finite-dimensional span of $1,j_1,\ldots,j_N$. This contradicts the infinite dimension of $J$, proving the assertion.
\end{proof}
